% interface=en output=pdftex

\environment mstyle

\TitlePage{514}{0}{\TeX UTIL\\explained}{}{1}

\subject {Introduction}

While typesetting a document, tables of contents, references and index
entries are often to be included. The items of a table of contents is
collected during the typesetting process, and therefore an additional
typesetting pass is needed to get the right one at the front of the document.
In a similar way, index entries are collected and need to be sorted before
inclusion. And, forward references can only be resolved when known. 

When processing a text, \CONTEXT\ saves table and index entries, cross
references, and some more information in a so called utility file. After a
successful run, this file is converted into a format suitable for the next
run. This conversion is done by \TEXUTIL, a \PERL\ script that comes with the
\CONTEXT\ distribution. 

As a bonus, \TEXUTIL\ can prepare documentation files, analyze graphics and
filter log files. These modes of operation will be discussed later. I will
not bother you with technical details on the file formats used. These can be
found in the \PERL\ file itself. 

To fill you in on the backgrounds of the \TEXUTIL\ logo as shown on this
cover: many lines of data produced by \CONTEXT\ are read in, interpreted,
converted, sorted and or combined, and written back to file in a similar or
adapted form: \TEXUTIL\ is dealing with transformations. 

{\em A next version will integrate (simple) conversion from \SGML\ to \TEX\ 
as well as sorting of mapped encodings.}

\subject {Running} 

\switch references   

When you use \TEXEXEC, running \TEXUTIL\ is taken care of automatically.
Otherwise, you need to call \TEXUTIL\ manually: 

\starttyping
texutil  --references  somefile
\stoptyping

In this case, \TEXUTIL\ loads the file \type {somefile.tui} sorts and
converts the content, and writes a file called \type {somefile.tuo}.
Additional switches may be passed; these will be discussed later. 

\switch documents   

The \CONTEXT\ related source files often contain documentation. The
documentation lines can be recognized by \type {%D} for text lines, \type
{%M} for macro code needed for the documentation, and \type {%S B} and \type
{%S E} for parts that do not make it into the documentation. In \PERL\
scripts, the~\type {%} is replaced by a~\type {#}. So, when you take a look
at source files, you now know what these symbols stand for. 

\starttyping
texutil  --documents  somefile
\stoptyping

This command converts a \TEX\ file with suffix \type {tex} into a file with
suffix \type {ted}. A related switch \type {--sources} does the opposite: it
removes the documentation lines and generates a source only file with suffix
\type {tes}. More information can be found in \type {texutil.pl}. 

\switch figures   

Although \PDFTEX\ can include graphics data, traditional \TEX\ depends on the
\DVI\ post||processor. Whatever method is used, \TEX\ has to reserve space
and therefore has to now the dimensions of graphics. Determining the
dimensions of for instance \EPS\ graphics can be done by \TEX\ itself, but
bitmap formats need to be analyzed by some other program, and that is where
\TEXUTIL\ comes into view. 

\starttyping
texutil  --figures 
\stoptyping 

When told to, \TEXUTIL\ starts searching for graphic files and collects some
information on them in a file called \type {texutil.tuf}. When a graphic is
included, and \CONTEXT\ cannot sort out the dimensions itself, it will search
for data in this file. When needed and enabled, \CONTEXT\ will call \TEXUTIL\
to provide the information. 

\starttyping
texutil  --figures  *.png  *.tif 
\stoptyping 

As demonstrated here, you can limit the inventory. This is useful when you 
want \CONTEXT\ to typeset a directory of figures. This kind of \TEXEXEC\ 
features are not that special; they just take care of calling the right 
\CONTEXT\ macro. 

\switch logfile   

\TEXUTIL\ is capable of concerns parsing the log file as produced by
\CONTEXT. The command: 

\starttyping
texutil  --logfile  somefile 
\stoptyping 

will report information on unknown references and overfull boxes: both being
important issues to resolve when typesetting a file. 

\switch purge  

Although \CONTEXT\ uses only one utility file, many scratch files can scatter
your disk after a run: buffers, \METAPOST\ graphics as well as files that 
communicate between \CONTEXT\ and \METAPOST, log files, etc. 

\starttyping
texutil  --purge 
\stoptyping 

Afterwards, \TEXUTIL\ reports the number of bytes regained. Optionally, you
can pass a file pattern, thereby limiting the cleansing. 

\subject {Switches} 

Before we will explain the mode dependant switched, we will introduce some
general ones. 

\switch help 

This switch results in some help information on the switches to use. 

\switch silent

This switch suppresses all messages to the screen. It enables the use of
\TEXUTIL\ in ways like: 

\starttyping
texutil  --figures  --silent --verbose somefile.png 
\stoptyping 

This reports something like:  

\starttyping
n=somefile.png t=png x=0bp y=0bp w=443bp h=591bp
\stoptyping 

This information can be piped into another program. 

\switch verbose

Sometimes there is more to report than useful at first sight. This switch
will therefore lead to more screen output. 

\switch interface=languagecode

If you want to let \TEXUTIL\ talk back in your own language, use this switch.
Currently \type {en}, \type {de} and \type {nl} are supported. 

\subject {Installation} 

Installation is not complicated. When you have \PERL\ running, \TEXUTIL\
should run too. On \MSWINDOWS\ you can copy \type {runperl.exe} into \type
{texutil.exe} to directly run \TEXUTIL. This program is part of the \FPTEX\
distribution. 

\subject {References}

The slightly misleading switch \type {--references} concerns more than
handling cross references. Actually these, as well as list entries and two
pass data, are simply passed on to the output file. Index entries and sorted
list entries however are sorted and redundant index entries are removed. 

\switch ij 

Sort \type{IJ} as \type {Y}. This is a typical dutch oriented switch. 

\switch high 

This switch converts high standard \MSDOS\ \ASCII\ values into \TEX\
commands. This switch is provided for backward compatibility reasons. 

\switch quotes 

Normally quotes are used as boundary characters for strings. This switch
forces \TEXUTIL\ to handle quotes differently. 

\subject {Documents}   

Documentation is part of the source files. The \type {--documents} option
lets \TEXUTIL\ enter a filter mode. Because typesetting of documentation is
supported by \TEXEXEC, users seldom have to deal with this option directly. 

\switch type=suffix   

Normally \TEXUTIL\ will try to sort out the right filetype, but when more
instances of a file (with different suffices) are present, it makes sense to
tell \TEXUTIL\ explicitly which filetype is to be converted. 

\switch outputfile=filename

This file can be used to force \TEXUTIL\ in using a different name for the
resulting output. 

\switch sources 

By default, documentation files are generated, but when this switch is
passed, source code only files are produced. 

\switch setups   

There used to be a time when syntax definitions were part of the source
files. Currently these are collected in one file, but maybe some day they
will make it to the source files again. This option collects all definitions
in a file \type {texutil.tus}. 

\switch infos   

The first few years of \CONTEXT\ all macro definitions were collected into a
few files. These files also contained explanations of the commands, that
could be used by help systems, as implemented in \TEXEDIT. While documenting
the sources, this information is removed, but the feature is still there. The
help information is recognized by \type {%I} directives and is collected in
files with the suffix \type {tud}. 

\switch templates 

This option is related to the previous one and concerns editor templates.
Again we're dealing with a features that might become obsolete. 

\subject {Figures} 

As mentioned, \TEXUTIL\ can analyze graphics (\PNG, \JPG, \TIF, \EPS, \PDF\
and \MPS). The next few switches deal with a different topic: conversion from
\EPS\ to \PDF. 

\switch epspage 

In preparation for conversion to \PDF, \EPS\ files should be prepared in such
a way that the bounding box is similar to the paper dimensions. This option
prepares the files and also removes some redundant (rubbish) information. To
prevent repetitive adaption the files are marked as being adapted. 
 
\switch epstopdf 

When issued, this option calls \GS\ to convert the \EPS\ graphic into a \PDF\
one. 

\subject {Logfile}

The faster machines get, the less obvious are messages. Filtering the log
files therefore makes sense. The \type {--logfile} option serves this
purpose. 

\switch box 

Issuing this switch lets \TEXUTIL\ report overfull boxes, both horizontal and
vertical. ALternatively you can use \type {--vbox} and \type {--hbox}. 

\switch criterium=number 

When reporting overfull boxes, small deviations, in the order of magnitude of
a few points, are seldom disturbing. By setting a criterium \TEXUTIL\ will
only report those overfull boxes that exceed the criterium. The number
represents points. 

\switch unknown 

A missing reference or an unknown figure is not what we expect to be
reflected in the final document. This option makes sure we don't miss them 

\ColofonPage

\stoptext
